% !TeX root = ../教材题库_全部条目_转写工作稿.tex
% Source: XJ-B1-C02-EX-P047-Y0578
% 第2章 / 2.2 从函数观点看一元二次方程 / 教材 p.47 / PDF 第54页
% OCR未校对：公式与文字均可能有误；未完成数学TeX转写；禁止直接发给学生
% 原题及相邻原解截图（核对用）：../题目原图/XJ-B1-C02-EX-P047-Y0578-01.webp
% 本文件只在转写完成后作为题干使用；OCR 条目尚不是可靠题干。
\begin{bankitem}{2.2 例题 例题2}{第2章 / 2.2 从函数观点看一元二次方程 / 教材 p.47 / PDF 第54页}{XJ-B1-C02-EX-P047-Y0578}{OCR未校对：公式与文字均可能有误；未完成数学TeX转写；禁止直接发给学生}
\begin{ocrdraft}
[教材p.47 / PDF第54页]
\par
00已知二次函数3，一十十c的图象与3，轴交于点A（0，
\par
一3），与工轴的
\par
两个交点的横坐标的平方和为巧，求该二次函数的表达式．
\par
解由二次函数的图象与3，轴交于点，A（0，一3）知，
\par
设二次函数的图象与轴交点的横坐标为，。，
\par
2《是
\par
一兀二次方程
\par
十一3一0
\par
的两个根，由根与系数的关系知
\par
XI．+一工2
\par
一兀二次方程a工2+
\par
所以
\par
十==()I十2）2
\par
+c一0的根与系数的关
\par
一b2十6一巧，
\par
系为．
\par
解得乃一\ensuremath{\pm}3．
\par
．丆1\ocrglyph{00AF}f一工2
\par
故所求二次函数的表达式为一十3一3或3，
\par
3x一3．
\end{ocrdraft}
\end{bankitem}
